% ===========================================================================
% Related Work -- extension covering the seven recent papers.
% Drop-in subsections plus a comparison table. Keys match new_refs.bib.
%
% All seven were read in full. Every number quoted below is from the paper
% itself, not from an abstract or an index entry.
%
% Organised by what each method's learned component decides, because that is
% the axis on which the seven separate most cleanly, and because both of the
% lineages among them fall inside a single group each -- so how the newer
% papers build on the older ones is told within the groups rather than
% needing a structure of its own.
% ===========================================================================


\subsection*{Deciding when to stimulate}

One line of recent work leaves frequency and amplitude fixed and learns only
whether to deliver stimulation at each step. Mehregan et
al.~\cite{mehregan2024enhancing} discretise the state and action spaces of an
actor--critic controller and reward reduced beta-band power while holding
stimulation frequency low, then add post-training and quantisation-aware
training to cut inference cost; they validate in a rodent model as well as in
simulation, which is the only in-vivo evidence among the seven. Ravivarapu et
al.~\cite{ravivarapu2025sea} continue that line directly, adopting the same
group's network model but targeting sample efficiency rather than model size.
A predictive reward model estimates outcomes so the policy need not query the
environment at every update, Gumbel-Softmax sampling handles exploration in
the binary action space, and FP16 quantisation reduces the model from 65\,MB
to 33\,MB. Their reported baseline is DDPG.

The two divide the same problem between them: one compresses a trained
policy, the other reduces the interaction needed to train it. They differ in
how stimulation cost is treated. Mehregan et al.\ carry it as a second reward
term, penalising stimulation frequency alongside beta power, whereas the
reward in~\cite{ravivarapu2025sea} is purely quadratic in the deviation of
mean beta power from a fixed threshold, with sparse stimulation emerging from
that form rather than being asked for.


\subsection*{Deciding the stimulation parameters}

Nguyen et al.~\cite{nguyen2026neuromorphic} tune the stimulus itself rather
than its timing. A deep spiking Q-network reads a biophysical cortico-basal
ganglia-thalamic model and emits nine discrete actions, one of
$\{-1, 0, +1\}$ for each of frequency, pulse width and amplitude, within
clinical bounds; the reward combines alpha--beta suppression with the physical
charge delivered. The policy is compressed by sparsity-constrained knowledge
distillation and deployed on neuromorphic hardware, where the authors report
45.2\% suppression, an 80.0\% reduction in cumulative stimulation charge
against continuous DBS, 0.52\,mW inference power, and a 28.1$\times$ reduction
in energy per inference against an equivalent conventional network on an
embedded GPU.

The distinguishing feature is that its action space is relative rather than
absolute. Each step nudges one parameter up or down instead of selecting a
setting, so the policy encodes a trajectory through the parameter space rather
than a choice within it, and the reachable settings at any moment depend on
where the previous steps left it.


\subsection*{Deciding which contact}

A separate line treats the electrode contact, not the waveform, as the
quantity to optimise, and comes from stimulation of the ventral
capsule/ventral striatum for psychiatric indications rather than from
Parkinson's disease. Nagrale et al.~\cite{nagrale2023bayesian} cast contact
selection as a multi-armed bandit, each contact an arm with an unknown
behavioural reward distribution, driven by UCB1 over reaction times on the
Multi-Source Interference Task, with amplitude, pulse width and frequency held
fixed. Their simulation environment is fitted to recorded data from six
participants.

Nagrale and Widge~\cite{nagrale2026arena} revise that design in three ways and
report what each is worth. Removing the intermediate state-space sensor model
and feeding the bandit raw reaction times lowers the contact ranking failure
rate from 31.9\% to 11.9\% and avoids a variance inflation of 51.7\% at
baseline and 148.0\% in conflict states. Replacing discrete daily resets with
one continuous history drives mean trial-level regret down steadily over
4{,}200 trials. And a rolling arena that keeps four of eight contacts active,
benches the weakest and rolls in an untested one scales to the full array
without sacrificing accuracy. The comparison spans UCB1, a variance-aware UCB,
Block-Variance Bayes-UCB and Thompson Sampling variants against greedy,
$\varepsilon$-greedy and brute-force search.

The pair is the clearest case here of a method correcting itself: the later
paper's central finding is that an estimator placed between the measurement
and the learner cost more in added variance than it recovered. Contact
selection and parameter tuning also remain complementary rather than competing
problems, one choosing where to stimulate and the other what to deliver,
and~\cite{nagrale2026arena} lists joint optimisation of amplitude, pulse width
and frequency as future work.


\subsection*{Tuning a controller rather than the stimulus}

Mai et al.~\cite{mai2026personalized} keep an explicit controller and use
learning only to set its parameters. An unknown-input observer is paired with
nonlinear state-error feedback, and a bounded soft actor--critic policy adjusts
two observer gains and two feedback gains every 20\,ms; the reward balances
tracking, smoothness and energy efficiency, and the entropy term is credited
with robustness under non-stationary coupling and cortical drive. On a delayed
subthalamic--globus pallidus externa firing-rate model whose dominant
oscillation sits near 20\,Hz, they report a mean RMSE of 0.4143 spikes/s
against 0.4816 for conventional active disturbance rejection control and
1.5000 for sliding mode control.

What is learned here is how to regulate, not what to deliver: the action space
is controller gains, and the stimulus is whatever the tuned controller
produces. The search space is correspondingly small, four gains rather than a
parameter grid or a pulse sequence. Note also that the work is titled as
personalised but the personalisation reported is robustness; the policy is
trained under randomised coupling and disturbance multipliers rather than
calibrated per individual, and no patient data or severity grading enters.


\subsection*{Suppressing synchrony in oscillator ensembles}

Krylov et al.~\cite{krylov2020framework} evaluate on ensembles of abstract
oscillators rather than a biophysical circuit, and take synchrony itself as
the pathology. Their environment spans Bonhoeffer--van der Pol,
Hindmarsh--Rose and Rulkov neuron models, and they demonstrate suppression
across three signalling regimes using multiple PPO agents on ensembles of
$N = 1000$ neurons. Their reward,
$R_t = -(X(t) - \langle X \rangle_t)^2 - \beta |A(t)|$ with $\beta = 2$,
penalises both residual synchrony and the amplitude of the applied action, and
they note that the action space can be constrained so an agent's outputs
remain harmless to the patient.

This is the one contribution among the seven that is an environment rather
than a controller, and its endpoint differs from every other: synchrony across
an ensemble, rather than a band power or a behavioural response. The
abstraction cuts both ways, giving a setting cheap enough to compare many
agents in but one whose relationship to a measured clinical biomarker is not
established.


% ===========================================================================
% Comparison table
% ===========================================================================

\begin{table}[t]
\centering
\caption{The seven recent works, grouped by what the learned component
decides. ``Stim.\ cost'' marks an explicit term for stimulation cost in the
objective, as opposed to sparse stimulation emerging from the form of a
biomarker-only reward.}
\label{tab:recent-related}
\small
\setlength{\tabcolsep}{4pt}
\begin{tabular}{@{}lllll@{}}
\toprule
Work & Decides & Environment & Endpoint & Stim.\ cost \\
\midrule
\multicolumn{5}{@{}l}{\emph{When to stimulate}}\\
\cite{mehregan2024enhancing} & Pulse on/off     & CBGT + rodent       & Beta power    & yes \\
\cite{ravivarapu2025sea}     & Pulse on/off     & CBGT (HH neurons)   & Beta power    & no  \\[2pt]
\multicolumn{5}{@{}l}{\emph{Stimulation parameters}}\\
\cite{nguyen2026neuromorphic}& Freq., PW, amp.\ & CBGT                & Alpha--beta   & yes \\[2pt]
\multicolumn{5}{@{}l}{\emph{Which contact}}\\
\cite{nagrale2023bayesian}   & Contact          & Fitted RT generator & Reaction time & no  \\
\cite{nagrale2026arena}      & Contact          & Fitted RT generator & Reaction time & no  \\[2pt]
\multicolumn{5}{@{}l}{\emph{Controller gains}}\\
\cite{mai2026personalized}   & Observer gains   & Delayed STN--GPe    & Firing rate   & yes \\[2pt]
\multicolumn{5}{@{}l}{\emph{Ensemble synchrony}}\\
\cite{krylov2020framework}   & Action amplitude & Oscillator ensemble & Synchrony     & yes \\
\midrule
This work                    & Freq.\ \& amp.\  & BGT                 & Beta power    & yes \\
\bottomrule
\end{tabular}
\end{table}

% Two notes on reading the table.
%
% The endpoints are not mutually comparable. Beta power, reaction time,
% firing-rate error and ensemble synchrony are different quantities, each
% measured against a baseline of that paper's own choosing, so the reported
% percentages cannot be ranked against one another. The table is grouped
% rather than ordered for that reason.
%
% The stimulation-cost column separates less than might be expected: five of
% the eight rows carry an explicit term. An energy-aware objective is now
% common rather than distinguishing, and the column is kept because that
% itself is worth showing.
